\section{Evaluation and Conformance Evidence}

Jankurai's evaluation claim in this edition is intentionally conformance-oriented: can the repository emit reproducible artifacts that make a merge decision auditable? This is not a broad empirical benchmark, a defect study, or a claim of detector completeness. It is evidence that the reference implementation can bind fixture inputs to observed audit and witness decisions.

\subsection{Current Implementation Evidence}

The \texttt{jankurai} CLI audits this repository and emits JSON and Markdown reports. It validates rule coverage from \texttt{agent/vibe-coverage.toml}, tracks generated zones, owner routes, proof lanes, score history, repair queues, and merge witnesses, and builds this paper from source. The paper, schemas, agent maps, generated-zone manifest, audit reports, conformance results, and generated coverage outputs are repository-local artifacts rather than hosted service state.

\begin{table}[t]
\caption{Current self-audit evidence for this paper cut.}
\label{tab:self-audit-evidence}
\centering
\scriptsize
\begin{tabularx}{\columnwidth}{L{0.36\columnwidth} Y}
\toprule
\JKTableHeader
\textbf{Artifact or count} & \textbf{Current value} \\
\midrule
Audit JSON/Markdown & \texttt{agent/repo-score.json}, \texttt{agent/repo-score.md}. \\
Score / findings & Current audit lane output under \texttt{agent/}. \\
Vibe coverage & Source rows validated from \texttt{agent/vibe-coverage.toml}. \\
Conformance run & Schema-valid JSON and Markdown under \texttt{target/jankurai/}. \\
Generated paper table & Regenerated by \texttt{just conformance}. \\
\bottomrule
\end{tabularx}
\end{table}

Current evidence remains bounded. Some detectors are deterministic, such as secret-like content, generated-zone mutation, direct DB markers, destructive SQL heuristics, unsafe input-boundary markers, and dead-language markers. Other rows are policy-backed or advisory because they require project-specific fixtures, organizational proof, or semantic tests that cannot be inferred from repository shape alone.

\subsection{Observed Seed Conformance Suite}

The seed conformance suite under \texttt{conformance/} is small by design. Each fixture carries a \texttt{jankurai-fixture.toml} manifest naming changed paths, expected audit decision, expected witness decision, expected rules, expected missing evidence, and optional proof receipts. The release lane is:

\begin{lstlisting}[style=codebox]
cargo run -p jankurai -- conformance run \
  --fixtures conformance/fixtures \
  --expected conformance/expected \
  --out target/jankurai/conformance-results.json \
  --md target/jankurai/conformance-results.md \
  --tex paper/tex/generated/conformance_results_table.tex
\end{lstlisting}

The runner executes the existing repository audit and reuses the merge-witness builder for each fixture. It then compares expected and observed audit decisions, expected and observed witness decisions, expected and observed stable rules, and expected missing-evidence text. The emitted report is schema-valid against \texttt{schemas/conformance-results.schema.json}; \texttt{just conformance} also runs the focused Rust tests that validate the report shape and generated TeX smoke output. Kubernetes conformance is a useful precedent because certification rests on repeatable tests and submitted evidence rather than branding alone \cite{kubernetesConformance2026}.

% Generated by: jankurai 1.5.1
% Source: conformance/fixtures; conformance/expected
% Command: cargo run -p jankurai -- conformance run --fixtures conformance/fixtures --expected conformance/expected --out target/jankurai/conformance-results.json --md target/jankurai/conformance-results.md --tex paper/tex/generated/conformance_results_table.tex
% DO NOT EDIT BY HAND.
\begin{table*}[t]
\caption{Observed conformance fixture results.}
\label{tab:conformance-results}
\scriptsize
\begin{tabularx}{\textwidth}{@{}L{0.22\textwidth} L{0.27\textwidth} L{0.12\textwidth} Y r@{}}
\toprule
Fixture class & Expected $\rightarrow$ observed & Rule & Result & ms\\
\midrule
% Fixture id: authz-isolation-fail
authz isolation & audit block$\rightarrow$block; witness block$\rightarrow$block & \texttt{HLT-022} & pass & 183\\
% Fixture id: destructive-migration-fail
destructive migration & audit block$\rightarrow$block; witness block$\rightarrow$block & \texttt{HLT-021} & pass & 130\\
% Fixture id: generated-zone-mutation-fail
generated zone mutation & audit block$\rightarrow$block; witness block$\rightarrow$block & \texttt{HLT-002} & pass & 131\\
% Fixture id: hl3-pass-minimal
hl3 pass minimal & audit pass$\rightarrow$pass; witness pass$\rightarrow$pass & \texttt{none} & pass & 135\\
% Fixture id: input-boundary-xss-fail
input boundary xss & audit block$\rightarrow$block; witness block$\rightarrow$block & \texttt{HLT-023} & pass & 153\\
% Fixture id: overbroad-agency-fail
overbroad agency & audit block$\rightarrow$block; witness block$\rightarrow$block & \texttt{HLT-012} & pass & 119\\
% Fixture id: ownerless-path-fail
ownerless path & audit block$\rightarrow$block; witness block$\rightarrow$block & \texttt{HLT-003} & pass & 123\\
% Fixture id: rendered-ux-gap-fail
rendered ux gap & audit block$\rightarrow$block; witness block$\rightarrow$block & \texttt{HLT-013} & pass & 143\\
% Fixture id: secret-sprawl-fail
secret sprawl & audit block$\rightarrow$block; witness block$\rightarrow$block & \texttt{HLT-010} & pass & 139\\
% Fixture id: unmapped-proof-fail
unmapped proof & audit block$\rightarrow$block; witness block$\rightarrow$block & \texttt{HLT-004} & pass & 130\\
\bottomrule
\end{tabularx}
\end{table*}


The current run observes 10 fixture matches out of 10: one passing HL3-shaped fixture and nine blocking fixtures. Covered failure classes include generated mutation, ownerless path, unmapped proof, secret sprawl, overbroad agency, rendered UX gap, destructive migration, authorization isolation, and unsafe input boundary. This strengthens the paper's conformance claim relative to inventory-only evidence: the release artifact now records expected-vs-observed decisions, not merely the presence of fixture files.

\begin{table}[t]
\caption{Seed fixture classes and expected signals.}
\label{tab:conformance-suite}
\centering
\scriptsize
\begin{tabularx}{\columnwidth}{L{0.28\columnwidth} Y}
\toprule
\JKTableHeader
\textbf{Fixture class} & \textbf{Expected signal} \\
\midrule
HL3 pass & Known-good shape with current proof receipt; audit and witness decision \texttt{pass}. \\
Generated drift & \texttt{HLT-002} blocks generated output changed without source proof. \\
Ownerless path & \texttt{HLT-003} blocks an unmapped changed path. \\
Unmapped proof & \texttt{HLT-004} blocks a path without a proof route. \\
Secret sprawl & \texttt{HLT-010} blocks secret-like fixture content. \\
Overbroad agency & \texttt{HLT-012} blocks tool or agent authority broader than the lane. \\
Rendered UX gap & \texttt{HLT-013} requires rendered evidence for changed UI. \\
Destructive migration & \texttt{HLT-021} blocks destructive SQL without safety evidence. \\
Authz isolation & \texttt{HLT-022} requires owner/non-owner negative proof. \\
Input boundary & \texttt{HLT-023} requires deterministic negative proof for unsafe input sinks. \\
\bottomrule
\end{tabularx}
\end{table}

\textbf{Miniature walkthrough.} In \path{generated-zone-mutation-fail}, the bad state is a generated output changed without its declared source contract and regeneration command. The expected finding is \path{HLT-002-GENERATED-MUTATION}. The witness blocks because a generated-zone touch requires source-regeneration proof. The repair queue should name the generated path, source path, regeneration command, owner route, and expected proof lane.

\subsection{Research Metrics}

\begin{table*}[t]
\caption{Metrics for measuring vibe-artifact reduction and agent efficiency.}
\label{tab:jankurai-metrics}
\centering
\small
\begin{tabularx}{\textwidth}{L{0.24\textwidth} Y}
\toprule
\JKTableHeader
\textbf{Metric} & \textbf{Meaning} \\
\midrule
Vibe artifact density & Open HLT findings per changed path, pull request, KLOC, or module. \\
Proof coverage ratio & Required proof lanes with current valid receipts divided by required lanes. \\
Owner route coverage & Changed paths matched to owner cells divided by all changed paths. \\
Generated-zone drift count & Hand-edited or stale generated outputs. \\
Repair half-life & Median time from finding to accepted repair receipt. \\
Context efficiency & Useful proof or retrieval success per context-pack token budget. \\
Witness completeness & Required witness fields present, current, and digest-bound. \\
Alignment drift & Delta from accepted conformance baseline. \\
Proof-selection accuracy & Required proof lanes selected without overbroad or missing lanes. \\
\bottomrule
\end{tabularx}
\end{table*}

The next empirical step is to measure whether the standard improves merge outcomes. Useful metrics include missing-evidence detection, false positives and false negatives, owner-route accuracy, proof-selection accuracy, review time, time-to-repair, CI runtime, regression prevention, security posture change, and token/context reduction. SetupBench and ContextBench show why environment bootstrap and context retrieval deserve direct measurement in agent workflows \cite{setupBench2025,contextBench2026}; Jankurai adds merge-evidence quality as the repository-local outcome to measure.
